\section{基函数变换}

一个常见的解决方案是将这样的单元与其邻居聚类形成一个聚合多面体，确保 $\mathbf{X_a}- \mathbf{X}_b$ 不会太小。聚合后，必须在这个新的多面体上定义基函数。虽然已有成熟的方法如虚单元法~\cite{MENGOLINI2019995}和调和基函数~\cite{Dasgupta2003}可用于多面体基函数，但它们在位移场的一致性和多材料单元局部区域（patch）应用方面存在挑战，并且线性虚单元法在多面体边界形状复杂时的精度也不高，而高阶虚单元法的积分则较为困难。

因此，我们提出了一种基于算子感知的局部基空间重构与隔离方法，利用现有的基函数和自适应 h-细化层次结构。\Cref{fig:refinement_adjoint_coarsening_demo} 展示了一个简单的两层示例：一个具有病态单元的网格（左上）被细化（上中），然后聚合成具有优化基函数的多面体（右上），其频谱分布与重新网格化的网格（左下）较为接近。实际上，病态单元及其周围环境具有多样性。本节介绍了一种用于构建网格层次结构的数值适应方法和相应的基函数优化策略。


\begin{figure}[h]
  \centering
  \includegraphics{figpdf/demo.pdf}
  \caption{\textbf{两层聚合示例。}一个具有明显病态单元的网格（左上）被细化（上中），然后聚合成具有优化基函数的多面体（右上），其频谱与重新网格化的网格（左下）接近。}
  \label{fig:refinement_adjoint_coarsening_demo}
\end{figure}


两个层级基函数之间的关系最初由细化过程中的几何插值定义：
\begin{equation}
  \varphi^k_i (X) = \sum_{j=1}^{n^k} \mathbf{C}_{ij}^k \varphi_{j}^{k+1}(X;z).
\end{equation}
其中 $z \in \mathcal{R}^3$ 是新添加顶点的位置，$\mathbf{C}^k$ 是一个 $3n^k \times 3 n^{k+1}$ 矩阵，其中 $\mathbf{C}_{ij} =
c_{ij}\mathbb{I}_{3\times3}$，标量系数 $c_{ij}^k$ 是重心坐标。需要注意的是，与 \cite{chen2019material} 中给定网格层次结构的系数 $c_{ij}^k$ 固定不同，在我们的方法中，由于 h-细化过程尚未完全指定，这些系数必须确定。鉴于我们的目标是局部操作而非全局优化，每一步 h-细化仅针对单个四面体单元及其直接邻居。选择该候选四面体单元的方法详见 \Cref{sec:selection}。假设当前层级 $\mathcal{M}^k$ 中的候选四面体单元已确定，我们的任务就变成了优化局部区域 $\mathcal{P}_i^H = \cup_j\mathcal{T}^k_j$ 的基函数，该局部区域（patch）由聚合多个四面体形成。

在 \cite{chen2019material, chen2018numerical} 中，他们通过最小化 Galerkin 刚度矩阵的迹 $\operatorname{tr}\left(\mathbf{C} \mathbf{K}^{k+1} \mathbf{C}^T\right)$ 来优化矩阵值基函数。考虑到
\begin{equation}
  \mathbf{K}^{k+1}_{\mathcal{P}_i}(z) = \int_{\Omega_{\mathcal{P}_i}}\varphi^{k+1}(X;z) \mathcal{L} \varphi^{k+1, T}(X;z)  dX,
\end{equation}
优化的材料适应基函数 $\mathbb{C}^k_{\mathcal{P}_i}$ 是最小化以下问题的解
\begin{equation}
  \begin{split}
    \min_{\mathbb{C}, z} &\trace{\mathbb{C}  \mathbf{K}^{k+1} (z) \mathbb{C}^T}\\
    \text{s.t.}  &\mathbb{C}^T \left(\vec{1}_{n^H}  \otimes \mathbf{I} \right) = \left(\vec{1}_{n^{h+}} \otimes \mathbf{I} \right), \quad \mathbb{C}^T \left( \left[ \mathbf{X}^H \right]_X \right) =  \left( \left[ \mathbf{X}^{h+} \right]_X \right)\\
    &\mathbb{C}_{i,j}^T = \delta_{ij} \mathbf{I}, \forall i,j \in \mathbf{v}^H, \quad \mathbb{C}^T \mathbf{h}_{ab}^H = \mathbf{h}_{ab}^{h+}
  \end{split}
  \label{eq:opt_basis}
\end{equation}
$\mathcal{P}_i$ 的下标表示域已被省略以简洁起见。所提供的约束确保刚体运动不变性和满足调和位移，这些使用 $\mathcal{P}_i$ 周围的一环邻域计算；更多细节可在~\cite{chen2018numerical, heFast2023} 中找到。虽然这种优化本质上耦合了细化和粗化，但其非线性和不连续性质，特别是当顶点放置在不同位置（例如边或面上）时网格拓扑和单元数量的变化，带来了优化挑战。为了简化这个复杂的优化问题，我们采用交替策略。

\paragraph{网格细分}
我们通过直接最小化细分后四面体块的迹来优化 $z$
\begin{equation}
    \min_{z}  \sum_{j, \Omega_{\mathcal{T}_j^{k+1}} \subset \Omega_{\mathcal{P}}} \trace{A^{k+1}_j (z)} \quad \text{s.t} \Vol{E^{h+}_i} > 0,\forall i \in E^{h+}.
  \label{eq:harmonic_refinement_trace}
\end{equation}
具体来说，我们分别为边细化求解六种情况，为面细化求解四种情况。对于面细化，两个四面体被分成六个，如 \Cref{fig:refinement_adjoint_coarsening_demo} 所示。对于边细化，受影响边的所有一环四面体被分成两个。然后我们按这些十种情况的最大特征值排序，选择产生最小值的配置。虽然可以使用涉及层级 $k+1$ 中额外新顶点的更复杂细化策略，我们将其留待未来工作。

\begin{figure}[h]
    \centering
    \includegraphics[width=0.4\linewidth]{figures/split-zjx/F4.0523a.png}
    \caption{\textbf{局部网格细分示意。}对于候选四面体及其局部邻域，我们分别考虑边细分与面细分等不同情形，并比较细分后局部四面体块的数值性质，从中选择更优的细分配置。}
    \label{fig:local-mesh-subdivision}
\end{figure}

一旦 $z$ 确定，\Cref{eq:opt_basis} 中的优化就变成了标准的二次规划（Quadratic Programming, QP）问题。由于我们总是将新顶点保留在局部区域（patch）内，边界顶点保持拉格朗日性质，避免了多面体边界处的任何不连续性问题。也可以通过遵循~\cite{chen2019material} 并将基函数投影到刚体运动约束流形上来优化基函数，尽管这与我们的方法相比差异很小。

\subsection{显式强连接与隐式强连接}

上述基于细分的基函数优化会在局部区域（patch）中引入一个新的离散定点，其位置由 \Cref{eq:harmonic_refinement_trace} 所定义的优化问题确定。若该优化点恰好对应于输入网格中的已有顶点，则称其为“显式强连接”。这意味着导致局部刚度异常增大的强耦合已经在顶点层面被明确暴露出来，后续可以直接围绕该顶点进行基函数重构与局部隔离。相反，若该优化点不对应于输入网格中的任何已有顶点，则称其为“隐式强连接”。这类强连接虽然尚未在原始顶点集合上直接显现，但其数值影响已经通过局部几何退化和刚度谱的异常增长表现出来。

从离散结构上看，基于细分的基函数优化本质上是将包含隐式强连接的病态单元局部细分为若干更小的四面体单元，使原本隐藏在单元内部的强耦合逐步显化为顶点之间可识别的显式强连接，如 \Cref{fig:strong-link-implicit-to-explicit} 所示。显式强连接对应的四面体通常具有两个距离极近的顶点；在后续的基函数优化中，我们将其中一部分局部高频模态从主计算空间中分离出来，从而隔离这类强耦合对整体离散系统的影响。与此同时，反映局部物理结构特性的插值关系则由优化得到的矩阵基函数系数给出。

\begin{figure}
    \centering
    \includegraphics[height=0.20\textwidth]{figures/subdivide/foot_at_center.png}
    \includegraphics[height=0.20\textwidth]{figures/subdivide/foot_at_edge.png}
    \includegraphics[height=0.20\textwidth]{figures/subdivide/like_quad.png}
    \caption{\textbf{隐式强连接经细分后转化为显式强连接。}}
    \label{fig:strong-link-implicit-to-explicit}
\end{figure}

从层次结构构建的角度看，强连接的处理可分为两个连续阶段：首先通过局部 h-细化将输入网格中的隐式强连接转化为细化网格中的显式强连接；随后在由 $\mathcal{M}^{h+}$ 到 $\mathcal{M}^{H}$ 的基函数重构过程中，将这些显式强连接对应的高频模态隔离到互补空间中。下面我们将这一过程统一纳入基于细化的层次结构构建框架中加以描述。

\subsection{基于细化的层次结构与强连接处理}

下面我们将上述过程统一描述为从 $\mathcal{M}^{h}$ 到 $\mathcal{M}^{h+}$ 再到 $\mathcal{M}^{H}$ 的两阶段流程。第一阶段通过局部 h-细化将输入网格中的隐式强连接显式化；第二阶段在细化网格上进行单元聚合与基函数优化，并对仍然存在的显式强连接进行隔离，从而得到最终用于计算的网格与函数空间。

\begin{figure}[htbp]
  \centering
  \begin{tikzpicture}[x=0.95cm,y=0.95cm]
    % unified styles (kept local to match repo convention)
    \definecolor{hierobox}{RGB}{245,246,250}
    \definecolor{hieroline}{RGB}{70,70,70}
    \definecolor{hierotext}{RGB}{40,40,40}

    % main boxes
    \node[
      draw=hieroline,
      fill=hierobox,
      thick,
      rounded corners=2pt,
      minimum width=2.75cm,
      minimum height=1.08cm,
      font=\normalsize
    ] (MH) at (0,0) {$\mathcal{M}^H$};

    \node[
      draw=hieroline,
      fill=hierobox,
      thick,
      rounded corners=2pt,
      minimum width=2.75cm,
      minimum height=1.08cm,
      font=\normalsize
    ] (Mh) at (4.9,0) {$\mathcal{M}^{h}$};

    \node[
      draw=hieroline,
      fill=hierobox,
      thick,
      rounded corners=2pt,
      minimum width=2.75cm,
      minimum height=1.08cm,
      font=\normalsize
    ] (Mhp) at (13.1,0) {$\mathcal{M}^{h+}$};

    % above labels
    \node[font=\small\bfseries, text=hierotext, align=center] at (0,1.15) {多面体四面体混合网格\\无强连接};
    \node[font=\small\bfseries, text=hierotext, align=center] at (4.9,1.15) {输入四面体网格\\隐式强连接 + 显式强连接};
    \node[font=\small\bfseries, text=hierotext, align=center] at (13.1,1.15) {细分后的四面体网格\\仅有显式强连接};

    % refinement arrow
    \draw[-{Stealth[length=2.8mm]}, thick, draw=hieroline] ([xshift=0.12cm]Mh.east) -- ([xshift=-0.12cm]Mhp.west)
      node[midway, above=0.30cm, align=center, font=\small, text=hierotext] {隐式强连接细分后\\转化为显式强连接};

    % basis optimization arrow from below
    \draw[-{Stealth[length=2.8mm]}, thick, draw=hieroline]
      (Mhp.south west) to[out=232, in=308]
      node[pos=0.52, below=1.18cm, align=center, font=\small, text=hierotext] {基函数优化与强连接隔离}
      (MH.south east);
  \end{tikzpicture}
  \caption{基于细化的层次结构示意图。输入网格 $\mathcal{M}^{h}$ 通过局部细分生成 $\mathcal{M}^{h+}$，使隐式强连接转化为显式强连接；随后通过基函数优化与强连接隔离构造最终计算网格 $\mathcal{M}^{H}$。}
  \label{fig:paper2-hierarchy-overview}
\end{figure}

给定一个四面体网格 $\mathcal{M}^h = \{\mathcal{V}^h, \mathcal{T}^h\}$，其包含 $n^h$ 个顶点和 $m^h$ 个四面体单元，对应于标准线性基函数 $\{\varphi_i^h\}_{i=1\dots n^h}$。我们的层次结构并非预先给定，而是在处理病态单元的过程中通过递归局部细化动态生成。具体而言，从当前层级网格 $\mathcal{M}^k = \{\mathcal{V}^{k}, \mathcal{T}^{k}\}$ 出发，一步局部细化将其转换为 $\mathcal{M}^{k+1} = \{\mathcal{V}^{k+1}, \mathcal{T}^{k+1}\}$，并定义 $\mathcal{M}^0 \equiv \mathcal{M}^h$。默认情况下，输入网格中的所有顶点在细化过程中均被保留，因此有
\begin{equation}
  \mathcal{V}^0 \subset \mathcal{V}^1 \subset \cdots \subset \mathcal{V}^q
\end{equation}
其中 $q$ 表示 h-细化的总迭代次数。我们将最终细化后的网格记为 $\mathcal{M}^{h+}\equiv\mathcal{M}^q$。

在此基础上，我们进一步构造最终计算所使用的网格 $\mathcal{M}^{H}=\{\mathcal{V}^{H},\mathcal{P}^{H}\}$，其中 $\mathcal{P}^{H}$ 由多面体单元组成。大多数多面体实际上对应于输入网格中的原始四面体，因此通常有 $|\mathcal{P}^{H}\cap \mathcal{T}^{h}| \cong |\mathcal{T}^{h}|$。顶点集 $\mathcal{V}^{H}$ 在一般情况下与 $\mathcal{V}^{h}$ 相同；只有在存在 \Cref{sec:strong_links} 中将要讨论的显式强连接时，我们才会将其中部分顶点对应的局部高频模态从主计算空间中移出，此时有 $\mathcal{V}^{H}\subseteq\mathcal{V}^{h}$。最终定义在 $\mathcal{M}^{H}$ 上的基函数记为 $\{\varphi_i^{H}\}_{i=1\dots n^{H}}$。

从空间分解的角度看，若暂不考虑对显式强连接的额外隔离，则细化网格 $\mathcal{M}^{h+}$ 上的函数空间可以分解为主计算空间与其互补空间之和，即
\begin{equation}
  \mathbb{V}^{h+} = \mathbb{V}^H \oplus \mathbb{W}^H.
  \label{eq:two-level-functiaonal-space}
\end{equation}
其中，$\mathbb{V}^{H}$ 由算子适应基函数 $\varphi^{H}=\mathbb{C}^{H}\varphi^{h}$ 张成，表示最终保留在主计算系统中的低频/主导模态；$\mathbb{W}^{H}$ 则表示与局部细化和基函数重构相对应的互补空间。换言之，$\mathcal{M}^{h}\rightarrow\mathcal{M}^{h+}$ 的过程负责将隐式强连接显式化，而 $\mathcal{M}^{h+}\rightarrow\mathcal{M}^{H}$ 的过程则负责在空间层面重构这些局部病态结构对整体系统的影响。

总之，我们首先通过若干步自适应 h-细化从 $\mathcal{M}^{h}$ 生成 $\mathcal{M}^{h+}$，使原本隐藏在病态单元内部的强耦合以显式强连接的形式暴露出来；随后通过单元聚合和算子适应基函数优化，从 $\mathcal{M}^{h+}$ 导出最终的计算网格 $\mathcal{M}^{H}$。这一过程的目标是在尽量保持自由度数量不变或仅小幅减少的前提下，降低使用 $\varphi^{H}$ 时全局刚度矩阵的最大特征值与条件数。

\subsubsection{显式强连接的隔离}\label{sec:strong_links}

经过上述局部细分后，输入网格中的隐式强连接已经在 $\mathcal{M}^{h+}$ 上转化为显式强连接。然而，对于这些仍然主导全局刚度矩阵最大特征值的显式强连接，仅靠局部细分和单元聚合仍不足以完全消除其影响，因而还需要在从 $\mathcal{M}^{h+}$ 到 $\mathcal{M}^{H}$ 的基函数重构过程中执行额外的隔离操作。

\begin{figure}
  \centering
  \includegraphics[width=0.7\linewidth]{plot_figures/Fig56/Fig56.0524.png}
  \caption{\textbf{显式强连接及其隔离效果。}左：一个典型的微结构模型，经 Tetgen 离散化为四面体网格后，其中部分单元包含一对距离极近的顶点，从而形成显式强连接。右上：$\mathbf{K}^H$ 与小波块 $\mathbf{B}^H$ 的频谱随迭代的变化。右下：是否进行强连接隔离时的收敛比率比较。}
  \label{fig:strong-link-demo}
\end{figure}

这类显式强连接对应于两个顶点之间极短的几何距离。在拉格朗日基函数下，这会在相应顶点之间产生异常增大的基函数梯度，使局部刚度矩阵出现强烈的高频刚性分量，如 \Cref{fig:strong-link-demo} 所示。与一般病态单元不同，这种由顶点极近引起的强耦合并不会随着单纯的 h-细化自动消失；相反，由于输入顶点被保留在整个细化过程中，其影响会持续传递到后续的层次结构中。

为此，我们采用“狄拉克细化（Dirac refinement）”策略~\cite{chen2019material, chen_multiscale_2021}，将强连接对中一个顶点对应的局部高频模态从主计算空间中移出，并并入互补空间。为了确定应优先处理哪一个顶点，我们仍使用 \Cref{eq:rowsum_max} 的准则，选择对最大行和贡献更大的顶点。这样一来，原本的两层空间分解便自然扩展为更细的三层结构：
\begin{equation}
  \mathbb{V}^{h+} = \mathbb{V}^H \oplus \mathbb{W}^H \oplus \mathbb{W}^{h}.
\end{equation}
其中，$\mathbb{V}^{H}$ 表示最终保留的主计算空间；$\mathbb{W}^{H}$ 由被识别并隔离的显式强连接顶点相关的基函数构成；$\mathbb{W}^{h}$ 则包含局部细化过程中新增顶点所对应的基函数。前者负责承接原始输入网格中由显式强连接引起的局部高频模态，后者负责承接细化过程中引入的局部补充自由度。对这些基函数的优化仍遵循 \Cref{eq:opt_basis} 中给出的局部优化过程，而小波函数的支撑区域则由相应强连接顶点的一环邻域确定。

在这种分解下，最终线性系统可以写成块对角形式：
\begin{equation}
  \begin{bmatrix}\mathbb{C}^H\\ \mathbf{W}^H \end{bmatrix} \mathbf{K}^{h+} \begin{bmatrix} \mathbb{C}^{H,T} & \mathbf{W}^{H, T}\end{bmatrix} = \begin{bmatrix} \mathbf{K}^H & \bold{0}\\ \bold{0} & \mathbf{B}^H \end{bmatrix}
  \label{eq:block-diagonal}
\end{equation}
其中通常有 $n^H \gg n^w$。这意味着由显式强连接主导的高频刚性分量被集中到较小的互补块 $\mathbf{B}^H$ 中，而主计算块 $\mathbf{K}^H$ 的最大特征值则显著降低。由于在大多数实验中 $\lambda_{\min}(\mathbf{K}^H)\cong \lambda_{\min}(\mathbf{K})$，因此这种分解通常能够显著改善条件数并加速迭代求解过程。

换言之，细分步骤的作用是将原本隐藏在病态几何内部的强耦合显式化，而强连接隔离步骤则进一步将这些已经显式化的局部高频模态从主计算空间中剥离出来。二者共同构成了从 $\mathcal{M}^{h}$ 到 $\mathcal{M}^{h+}$ 再到 $\mathcal{M}^{H}$ 的完整层次结构构建过程。

\subsection{候选选择}\label{sec:selection}

为了在上述两阶段流程中优先处理对全局刚度矩阵谱性质影响最大的局部区域，本节给出候选四面体单元的选择准则。其目标不是重新定义层次结构本身，而是为局部细分、单元聚合以及后续的强连接隔离提供一个统一的排序依据，从而在每一步都优先处理最值得改进的病态区域。为简单起见，我们首先考虑同质情形，即所有材料具有相同的系数。为了降低全局刚度矩阵 $\mathbf{K}^{H}$ 的最大特征值（为简洁起见，以下省略上标 $H$），我们采用两个互补的不等式作为排序基础：

第一个来自 Gershgorin 圆盘定理：$ \max_i K_{ii} \le \lambda_{\max} \left( K \right) \le \max_i (K_{ii} + \sum_{j\neq i} |K_{ij} |)$。
第二个来自 Shewchuk 的“什么是好的线性单元”~\cite{Shewchuk2002whatIsAGoodLinearFiniteElement}：$\max_e \lambda_{e,\max} \le \lambda_{\max}\left( K \right) \le C \max_e \lambda_{e,\max}$。我们利用这些上界来确定要聚合哪些单元。由于两个界都与单元刚度相关，我们在组装过程中存储这些值。

后一个界自然给出了单元排序准则：
\begin{equation}
  e = \arg \min_{e} \lambda_{e,\max}
  \label{eq:elewise_eigmax}
\end{equation}
前一个界也可以通过单元贡献分解为：$K_{ii} + \sum_{j \neq i} | K_{ij} |=\sum_e \left( K^{e}_{ii} + \sum_{j\neq i} |K^{e}_{ij}| \right)$。然后我们选择对最大行和贡献最大的单元：
\begin{equation}
  e = \arg \max_e ( K_{ii}^e + \sum_{j \neq i} | K_{ij}^e |) \quad \text{s.t.} \quad i = \arg \max_{i} (K_{ii} + \sum_{j \neq i} |K_{ij}| ).
\label{eq:rowsum_max}
\end{equation}
当单元是聚合产生的多面体时，我们递归地将任一排序函数应用于它，将其视为局部网格，以找到进一步细化的特定候选四面体。

尽管采用了这种层次结构构建与基函数优化策略，最大特征值的上界在单级细化后并不总能显著降低。设 $y^k$ 表示经过第 $k$ 层细化后 $\mathbf{K}$ 的上界。如果相对变化 $\frac{\tilde{y}^{k+1} - y^k}{y^k}$ 低于阈值 $t$（我们默认设为 0.1），则继续执行下一轮局部细分与基函数重构。
